\chapter*{Abstract}
\addcontentsline{toc}{chapter}{Abstract}
\begin{spacing}{1.0}

An optimizing compiler improves a program by applying a sequence of transformation passes to an intermediate representation of the code. The passes interact: each one can create, destroy, or hide opportunities for those that follow, so the order in which they run changes the quality of the generated code. Choosing that order is the \emph{phase-ordering problem}. It remains unsolved in general because the number of orderings grows factorially and the best order depends on the program, so production compilers settle it with fixed, hand-tuned pipelines that hide the effect of ordering.

This project studies the problem directly by building a complete, transparent compiler infrastructure in which every optimization pass is independent and freely reorderable. The compiler accepts a subset of C with integers, fixed-size arrays, globals, functions, conditionals, and loops, and consists of a hand-written lexer, a recursive-descent parser, a semantic analyzer, a three-address-code intermediate representation with 35 opcodes, an interpreter that executes the intermediate code to obtain dynamic instruction counts and validate output, and a RISC-V RV32IM assembly back end. Six classical optimizations were implemented as reorderable passes: constant folding, dead code elimination, common subexpression elimination, copy propagation, strength reduction, and algebraic simplification. A pass manager runs any ordering and enumerates all $6! = 720$ permutations together with the unoptimized baseline.

Fifteen kernels adapted from the MiBench embedded suite and the PolyBench/C numerical suite were compiled under every one of the 721 orderings, giving 10{,}815 compile-and-execute experiments, each measured for static code size, an estimated cycle count from an ARM Cortex-M class cost model, and dynamic instruction count, then aggregated with the geometric mean. The best ordering family reduces code size by 35.7\% (22.3\% to 60.0\% per program), estimated cycles by 20.5\%, and dynamic instructions by 23.2\%; the worst family reduces code size by only 1.3\%. The 720 orderings collapse into nine equivalence classes determined solely by the relative order of four passes, and a simple rule governs them: constant folding and copy propagation must run early and dead code elimination last. Every optimized program produced output identical to its baseline, and the improvement survives lowering to RISC-V assembly.
\end{spacing}

\vspace{0.5em}

\noindent\textbf{Keywords:} Compiler Optimization, Phase-Ordering Problem, Intermediate Representation, Optimization Passes, MiBench, PolyBench/C, RISC-V, Exhaustive Permutation Analysis, Compiler Education.

